\section{Аналитический раздел}

\subsection{Анализ существующих решений}
В сфере обеспечения отказоустойчивости распределенных систем существует
некоторое количество решений, являющихся, де-факто, стандартом в данной области,
например:
\begin{itemize}
    \item ITRS Geneos~\cite{ref:geneos};
    \item Prometheus~\cite{ref:prometheus} + Grafana~\cite{ref:grafana}
\end{itemize}

Для изучения актуальности работы было проведено сравнение существующих решений
по следующим критериям:
\begin{itemize}
    \item основное назначение;
    \item модель данных;
    \item метод выявления проблем;
    \item анализ влияния сбоев.
\end{itemize}

В результате проведенного сравнения была составлена таблица~\ref{tbl:analogs-comparison}.

\begin{center}
    \begin{tabularx}{\textwidth}{|X|X|X|}
        \caption{\label{tbl:analogs-comparison} Сравнение существующих решений} \\ \hline
        Критерий сравнения & ITRS Geneos & Prometheus + Grafana \\ \hline
        \endfirsthead

        \multicolumn{3}{r}{{\tablename\ \thetable{} -- продолжение}} \\ \hline
        Критерий сравнения & ITRS Geneos & Prometheus + Grafana \\ \hline
        \endhead

        \hline
        \endfoot
        \endlastfoot

        Основное назначение &
        Контроль состояния серверов и приложений в реальном времени. &
        Сбор и визуализация метрик, анализ трендов. \\ \hline

        Модель данных & Иерархическая & Временные ряды \\ \hline
        Метод выявления проблем &
        \textbf{Реактивный}. Сравнение текущих значений с пороговыми и правилами. &
        \textbf{Реактивный}. Мониторинг, отклонений, оповещения. \\ \hline

        Анализ влияния &
        Ручной / Правила. Требует сложной настройки корреляции событий. &
        Отсутствует. Показывает состояние метрик, но не моделирует зависимости сбоев. \\ \hline
    \end{tabularx}
\end{center}

\subsection{Анализ предметной области}
В результате проведения анализа было выявлено, что существующие решения позволяют проводить только
реактивный анализ и не дают возможности изучать узкие места системы на этапе проектирования.
В связи с этим было решено, что разрабатываемая система должна предоставлять возможность
проактивного анализа для выявления узких мест и архитектурных проблем на этапе
проектирования системы.

Предметная область проекта охватывает управление надежностью распределенных программных систем
и анализ топологии микросервисной архитектуры. Система предназначена для моделирования
IT-инфраструктуры,
что позволяет автоматизировать поиск единых точек отказа и прогнозировать распространение
каскадных сбоев.

Основными сущностями предметной области являются:
\begin{itemize}
    \item компонент, может представлять одну из следующих сущностей:
        \begin{itemize}
            \item микросервис;
            \item база данных;
            \item брокер сообщений;
            \item внешний API;
        \end{itemize}
    \item связь между компонентами, может иметь различный уровень критичности.
\end{itemize}

\subsection{Ролевая модель}
Пользователи проектируемого приложения могут иметь следующие роли:
\begin{itemize}
    \item незарегистрированный пользователь;
    \item разработчик;
    \item SRE;
    \item архитектор;
    \item администратор.
\end{itemize}

На рисунке~\ref{img:usecases.pdf} представлена диаграмма прецедентов для описанных ролей.
\jimg{usecases.pdf}{Диаграмма прецедентов}

\subsection{Требования и ограничения к разрабатываемой базе данных}
В результате анализа были выделены основные сущности предметной области:
компоненты, выполняющие какие-либо задачи, например: отправка писем по электронной почте,
прием платежей, получение внешних данных, хранение данных, отправка сообщений и так далее,
и связи между этими компонентами, которые показывают от каких компонентов зависит работа
других компонентов, например: компонент для приема платежей зависит от базы данных
платежей, компонент получения внешних данных зависит от какого-либо API и очереди
сообщений для передачи этих данных, и так далее. Такого рода сеть можно представить
в виде ориентированного графа, что позволяет наглядно отобразить компоненты
и их зависимостей.

Таким образом основным требованием к разрабатываемой базе данных является возможность
оптимизированно хранить графы и эффективно выполнять операции над ними.

Так как было решено использовать ролевую модель на уровне приложения, помимо БД для хранения
графов необходимо спроектировать БД для хранения пользователей системы. Основным требованием
к данной БД является структурированное и непротиворечивое хранение данных.

Вот переформулированный и сжатый вариант текста в формате LaTeX с выделением графовой
модели в разделе постреляционных баз данных:

\subsection{Анализ моделей баз данных}

База данных представляет собой самодокументированное собрание интегрированных записей, 
структурированную согласно определенной модели. Традиционно выделяют три этапа развития 
моделей данных: дореляционный, реляционный и постреляционный.

\subsubsection{Дореляционные модели}
К данной категории относятся иерархическая и сетевая модели.

Иерархическая модель строится по принципу древовидной структуры, где записи
связаны отношениями <<предок--потомок>>. Реализуется связь типа 1:N, при
этом каждый потомок имеет строго одного предка. Удаление родительской записи влечет
автоматическое удаление дочерних, что обеспечивает ссылочную целостность.

Сетевая модель расширяет иерархическую, позволяя реализовывать связи типа 
<<многие--ко--многим>>. В такой модели запись может иметь несколько предков, что
превращает структуру в граф. Однако сложность организации памяти и жесткость схемы
затрудняют внесение изменений в структуру данных.

Достоинством дореляционных моделей является высокая скорость доступа к данным. К
недостаткам относятся значительные затраты памяти на хранение указателей,
сложность модификации структуры и перегруженность логики приложения деталями
низкоуровневого доступа.

\subsubsection{Реляционная модель}
В реляционной модели данные представляются в виде совокупности взаимосвязанных таблиц, 
а манипулирование ими осуществляется с помощью языка SQL.

Основные понятия модели:
\begin{itemize}
    \item домен --- множество допустимых значений;
    \item атрибут --- именованная характеристика сущности (столбец);
    \item кортеж --- набор значений атрибутов, образующий одну запись (строка);
    \item первичный ключ --- уникальный идентификатор кортежа;
    \item внешний ключ --- средство установления связей между отношениями.
\end{itemize}

Преимуществами реляционной модели являются высокая согласованность данных, простота инструментария и
независимость прикладного ПО от физического способа хранения данных.
Недостатками же являются снижение производительности при росте объема таблиц и жесткие
ограничения на структуру.

\subsubsection{Постреляционные модели}
Данные модели призваны преодолеть ограничения SQL и реляционного подхода, обеспечивая
гибкость при работе со сложными структурами. В отличие от классических БД,
здесь допускается использование многозначных полей и отказ от атомарности данных.

В рамках постреляционного подхода выделяют:
\begin{itemize}
    \item объектно-ориентированные и многомерные модели, снимающие ограничения
        на структуру записей;
    \item графовую модель, где данные представляются в виде узлов и ребер, 
        обладающих собственными свойствами, такая система оптимальна для работы с
        высокосвязными данными где важна
        скорость обхода сложных графов связей.
\end{itemize}

Достоинствами постреляционных моделей являются эффективная обработка огромных объемов неструктурированной
информации и поддержка большого числа пользователей.
Основным недостатком является сложность обеспечения строгой целостности и
непротиворечивости данных по сравнению с реляционными моделями.
